\cotefile{157-3}
\scanpage{222}
\mspage{172}
% L1–9
\uline{NB} Le groupe $\pi_1(X,x)$ est défini via la théorie du groupoïde fondamental. On sait qu'il est en bijection canonique avec $\pi_0(\Omega^1(X,x))$, où $\Omega^{\deleted{\unsure{1}}}(X,x)=\underline{\mathrm{Ch}}(X;x,x)$ (ce qui permet d'avoir une loi de composition monoïdale sur $\Omega^1(X,x)$, donnant la loi de groupe sur le $\pi_0$), ou au choix $\Omega^{\deleted{\unsure{1}}}(X,x)=\mathrm{Ch}_{\infty}(X;x,x)$ (plus commode dans certaines situations, mais ne donnant pas \inserted{directement} la loi de composition du $\pi_0\,\unsure{{}^{1}\!\Omega|\Omega}(X,x)$).

% L10
Je vais poser, pour $n\in\mathbb N$

% L11
\[
\pi_n(X,x)=\pi_0(\underbrace{\Omega^n(X,x)}_{\deleted{\unsure{\Omega}}}), 
\]
% L12–17
$\Omega^n$ désignant l'itéré $n$\unsure{.}ième du foncteur $(X,x)\longrightarrow\Omega(X,x)$ de la cat. \inserted{$\mathrm{Cat}_{\bullet}$} des \deleted{catégories} objets ponctués de $\mathrm{Cat}$ dans elle-même. Je préfère \deleted{\unsure{Je suis}} ici travailler avec la \unsure{méthode|théorie|modèle} $\underline{\mathrm{Ch}}$ au lieu de $\underline{\mathrm{Ch}}_\infty$. Tôt ou tard il faudra donner un énoncé de comparaison\ldots

% L18
Je veux définir $W_n$ pour $n\geq 2$.

% L19–21
\uline{Théorème\inserted{\unsure{$\checkmark$}} 1} Soit $f:X\to Y$ dans $\mathrm{Cat}$, \deleted{\unsure{soit donc}} \inserted{et soit $n\in\mathbb N$, $n\geq2$,} \deleted{\unsure{$S$}} une partie de $\mathrm{Ob}\,X$ telle que \deleted{\unsure{la catégorie}} $S\hookrightarrow\mathrm{Ob}\,X\longrightarrow\pi_0(X)$ soit \inserted{épi.} les conditions suivantes sont équivalentes

% L21, left margin
% slanted
\margin{Sans doute \uline{faux}. On doit avoir $a\Rightarrow b$, mais $b\Rightarrow a$ sans doute faux déjà pour $n=2$ (\unsure{sauf hyp.|sauf cas} $X,Y$ \unsure{simplement connexes}) \unsure{énoncé rectifié plus|énoncé rectif. post}.}


% L22–23
a) \deleted{$f\in W_1$ (cf page 171 ci-dessus) et} $\pi_i(f,x)$ iso pour tout $x\in\mathrm{Ob}\,X$, et \unsure{ts|les} $i\in[0,n]$


% L24
a\unsure{$'$}) Comme a), mais avec $x\in S$,


% L25–27
b) On a
\[
\begin{aligned}
 H^0(\unsure{Y|X},F)&\xrightarrow{\sim} H^0(X,f^*F) &&\text{pour $\unsure{F}$ faisceau d'ens loc. const. sur $Y$}\\
 H^1(Y,F)&\simeq H^1(X,f^*F) &&\text{si de plus $F$ faisceau en gr.}\\
 H^{\unsure{i|0}}(Y,F)&\simeq H^i(X,f^*F) &&\text{si de plus $F$ faisceau abélien}\\
 &&&i\in[1,n],
\end{aligned}
\]
% archive sheet number: 222
